Physics-informed neural networks (PINNs) are often presented as mesh-free solvers for fluid flow, but the
combination of remedies that makes them train is rarely tested against independent references with its
failures reported. We implemented that combination (Fourier features, gated networks with weight factorisation,
exact boundary and divergence constraints, causal and gradient-norm loss weighting, adaptive resampling,
separable networks and a physics-informed DeepONet) in one JAX code base and ran it on six problems with
acceptance criteria fixed in advance. The two-dimensional Taylor--Green vortex was solved to a velocity error
of $4.05\times10^{-5}$ and the inverse cylinder-wake problem recovered both equation coefficients to better than
0.1\%. On the lid-driven cavity at $\Rey=1000$ the best network is 0.27\% from a converged finite-difference
solution of the same problem, yet misses a 2\% gate against the classical Ghia tables, which an exact solution
also misses. Imposing the regularised lid exactly, an impulsive inflow and gradient-norm weighting each produced
networks that violate mass or energy conservation, and the cylinder and three-dimensional Taylor--Green gates
were not met within the reduced budgets the available GPU time allowed. A physics-informed DeepONet trained over
$100\le\Rey\le1000$ predicts the cavity at an unseen $\Rey=400$ with a 10.1\% error, just above its 10\% gate. We
report all runs (30 GPU hours in total), their costs and the reasons for each failure.
